\section{Limitations and Future Work}
\label{sec:limitations}

Our evaluation acknowledges several limitations. TF-IDF captures lexical patterns but may miss semantic relationships that deep learning approaches could reveal. The CEAS-08 corpus represents 2008 spam characteristics; modern corpora may interact differently with LLM-generated content. Our fixed class imbalance represents one operational point; different imbalance ratios may exhibit different patterns. Our mid-tier models may not generalize to open-source LLMs, frontier models, or specialized approaches including fine-tuned models. Our cross-model detection focus leaves same-model scenarios unexplored. Spam detection findings may not generalize to other cybersecurity domains requiring different feature representations or threat models.

To address these limitations, future work will expand our evaluation in several directions. We plan to incorporate additional corpora including Nazario, Enron, and MeAJOR datasets, paired with deep learning classifiers such as Sentence-BERT models to capture semantic features beyond lexical patterns. Furthermore, we intend to integrate transformer-based architectures for comparative analysis, assessing their generalization capabilities across broader cybersecurity domains including social engineering detection and adversarial content generation. These extensions will provide a more comprehensive understanding of LLM detectability across diverse threat landscapes and architectural approaches.
